\section{Lesson 0: What Demonstration Meant}

Open a page of Aristotle or St.\ Thomas Aquinas, and you meet prose that sounds at once serene and formidable. A question is posed; objections follow, armed with quotations from Scripture, the Church Fathers, or the philosophers; a brief counterpoint intervenes; and a compact paragraph settles the issue using ``motion,'' ``act,'' ``potency,'' ``form,'' ``essence,'' and ``demonstration.'' Everything arrives in the same brisk declarative voice. Without scholastic training, it is hard to tell when the author is defining a word, reporting a received opinion, leaning on an earlier proof, or claiming to prove something about reality from scratch.

The first hurdle is the word ``demonstration'' itself. A modern reader usually hears ``demonstration'' as ``a logically valid proof'' or ``a convincing argument.'' Neither captures what Aristotle meant by \emph{apodeixis} or what Aquinas meant by \emph{demonstratio}. For both, calling an argument a demonstration was a narrower and far stricter claim. A valid argument whose premises are all true can still fail to be a demonstration; conversely, an argument that is not a demonstration can still be indispensable to philosophy.

Before taking up the six reading questions, we need to see what demonstration, dialectic, disputation, and causal proof meant in the texts themselves. We also fix one rule of historical honesty at the threshold: our six reading questions---\textbf{Claim}, \textbf{Meaning}, \textbf{Warrant}, \textbf{Burden}, \textbf{Consequence}, and \textbf{Denial}---together with the eight kinds of warrant and five kinds of denial introduced in later lessons, are modern reading devices. They make distinctions the tradition cared about visible to a reader who never passed through a medieval arts faculty. They are never to be attributed to Aristotle or Aquinas, and whenever a modern term of ours stands beside a classical one, the difference between them stays explicit.

\subsection{Scientific Knowledge and the Conditions of Demonstration}

Aristotle's account of demonstration stands in Book I of the \emph{Posterior Analytics}. Because chapter divisions vary across medieval commentaries and modern editions, the Bekker page and line numbers govern every citation. At \emph{Posterior Analytics} I.2, 71b8--16 (echoed at \emph{Nicomachean Ethics} VI.3, 1139b18--22), Aristotle states what scientific knowledge (\emph{episteme}, Latin \emph{scientia}) is. Do not hear modern experimental science in that word. For Aristotle, \emph{episteme} names a state of the knower: we know a thing without qualification when we know the cause on which the fact depends, know it as the cause of that fact and of no other, and know that the fact cannot be otherwise. Its object is necessary, not accidental or merely usual.

Aristotle then defines demonstration (\emph{apodeixis}) at 71b17--19: ``By demonstration I mean a syllogism productive of scientific knowledge'' (G.\ R.\ G.\ Mure's Oxford translation). Because demonstration produces knowledge of a necessary fact through its cause, formal validity is its cheapest part. The real weight sits in six conditions on demonstrative premises (\emph{Posterior Analytics} I.2, 71b19--72a8): they must be \textbf{true}, \textbf{primary}, \textbf{immediate}, \textbf{better known than the conclusion}, \textbf{prior to it}, and \textbf{causes of it}.

Every word in that list carries weight. Truth excludes valid inferences from falsehoods. Being \textbf{primary} and \textbf{immediate} means having no prior premise standing behind them through which they must be proved; because neither circular reasoning nor an endless regress of proofs yields knowledge, every demonstrative science rests on undemonstrated starting points (\emph{archai}, Latin \emph{principia}) (\emph{Posterior Analytics} I.2, 72a8--25; I.3, 72b5--73a20). Being \textbf{better known} and \textbf{prior} is explicitly divided at 71b29--72a5 into two senses: better known \emph{to us} (what stands nearest sense perception) and better known \emph{in nature} (what is more universal and foundational in reality). Demonstrative premises must be better known \emph{in nature}. Finally, the premises must be \textbf{causes} of the conclusion: the shared middle term (\emph{meson}) linking the premises is not merely a verbal connector, but the actual cause in virtue of which the predicate belongs to the subject (I.4, 73a34--b26; II.11, 94a20--95a9).

Aristotle tightens the standard even further. Demonstrative premises hold necessarily and \emph{per se}---the predicate belongs in the subject's definition, or the subject in the predicate's---barring demonstrations of chance conjunctions or incidental accidents (I.4, 73a34--b26; I.6, 74b5--75a17; I.30, 87b19--27). They reach the commensurate universal, belonging to every instance of the subject precisely as such (I.4, 73b27--74a3), and stay within one subject genus unless one science stands subordinate to a higher one, as optics does to geometry or harmonics to arithmetic (I.7, 75a38--b20; I.13, 78b34--79a16). Because first principles cannot be demonstrated, Aristotle explains in \emph{Posterior Analytics} II.19 (99b18--100b17; cf.\ \emph{Nicomachean Ethics} VI.6, 1140b31--1141a8) that the universal settles in the soul from sense perception through memory and experience by \emph{epagoge} (induction), and \emph{nous} (intellect) grasps the first premises---a famous chapter whose interpretation scholars still dispute.

This single section already changes how you read metaphysics. Confronted with a valid syllogism built from true premises, you can now ask whether its premises are immediate, necessary, prior in nature, and causal. If any condition fails, the argument may still be true and useful, yet it is not a demonstration in Aristotle's sense.

\subsection{Dialectic, Disputatio, and Reading the Scholastic Article}

Because Aristotelian demonstration sets so high a bar, much serious philosophical reasoning cannot be demonstrative. In the \emph{Topics}, Aristotle separates demonstration from \textbf{dialectic}. Whereas demonstration reasons from primary, necessary, causal premises, dialectic reasons from \emph{endoxa}---reputable opinions accepted by everyone, by most people, or by the wise, and among the wise by all, by most, or by the most noted philosophers (\emph{Topics} I.1, 100a18--b23).

Dialectic is not failed demonstration. It is a distinct art with work demonstration cannot do (\emph{Topics} I.2, 101a25--37). It tests received opinions against one another, exposes hidden contradictions, clears conceptual ground, and provides the path toward the first principles of all inquiries---because starting points cannot be demonstrated from anything prior and can only be examined and defended dialectically. Recognizing that Aristotle or Aquinas is arguing dialectically is accurate reading, not a rebuke.

That practice shaped medieval \emph{disputatio}, the structured university debate in which a thesis was tested against its ablest objections. Every article of Aquinas's \emph{Summa Theologiae} follows this layout, which tells a reader which sentences on the page are actually Aquinas's own:

\begin{itemize}
\item The \textbf{question} (\emph{Utrum...}, ``Whether...'') states a live disputed issue.
\item The \textbf{objections} (\emph{Videtur quod...}, ``It seems that...'') give the strongest case against Aquinas's position, often quoting Scripture, Church Fathers, and Aristotle. Sentences inside an objection are \emph{not} Aquinas's own assertions.
\item The \textbf{\emph{sed contra}} (``On the contrary'') cites a brief authoritative counterpoint opening the other side, not yet Aquinas's philosophical proof.
\item The \textbf{\emph{corpus}} (\emph{Respondeo dicendum}, ``I answer that'') states Aquinas's own determination and reasoning.
\item The \textbf{replies} (\emph{Ad primum...}) answer each objection one by one, conceding what is sound, distinguishing senses, and narrowing cited authorities.
\end{itemize}

Knowing this layout keeps a reader from lifting an authority out of an objection---or mistaking a brief \emph{sed contra} for a full proof---and reading an article backwards.

\subsection{Showing That and Showing Why: Quia and Propter Quid}

Even within demonstration itself, Aristotle divides knowing \emph{that} a fact holds (\emph{tou hoti}, Latin \emph{quia}) from knowing \emph{why} it holds (\emph{tou dioti}, Latin \emph{propter quid}). In \emph{Posterior Analytics} I.13 (78a22--79a16), the two diverge in four ways: when premises are not immediate; when two attributes reciprocate and the better-known effect stands as middle term instead of the cause; when the middle term is too remote (saying a wall does not breathe because it is not an animal, when many animals do not breathe either, 78b13--34); and when one science shows the fact while a superior science supplies the reason (78b34--79a16).

Aquinas adopts this pair explicitly in \emph{Summa Theologiae} I, q.\,2, a.\,2, \emph{corpus}. In the Leonine Latin text he writes: \emph{Duplex est demonstratio}---``Demonstration is twofold.'' One proceeds through the cause, is called \emph{propter quid}, and argues from what is prior without qualification (\emph{per priora simpliciter}). The other proceeds through the effect, is called \emph{demonstratio quia}, and argues from what is prior to us (\emph{per ea quae sunt priora quoad nos}). Whenever an effect is better known to us than its cause, we can demonstrate \emph{that} its proper cause exists, since no effect exists without its cause. (Note that the 1920 English Dominican translation renders \emph{propter quid} as ``\emph{a priori}'' and \emph{quia} as ``\emph{a posteriori}.'' Because modern philosophy uses \emph{a priori} to mean ``independent of experience,'' retaining Aquinas's Latin terms \emph{quia} and \emph{propter quid} preserves the exact distinction of causal direction.)

This distinction sharpens every reading of natural theology. Proving \emph{that} a cause must exist is a real achievement, yet it does not show \emph{why} the effect flows from the cause's essence. That is why Aquinas states in \emph{ST} I, q.\,2, a.\,2, \emph{corpus} and \emph{ad 2} that the Five Ways (\emph{quinque viae}, q.\,2, a.\,3---his word is \emph{viae}, ``ways,'' not ``Five Proofs'') proceed by demonstration \emph{quia}, from effects better known to us, taking what the name signifies in place of an essence we do not directly see.

\subsection{What Formation a Modern Reader Lacks---and How We Avoid Anachronism}

Why does Aquinas's prose still feel opaque even with these terms in hand? Because medieval arts training formed three habits modern education rarely teaches:

\begin{itemize}
\item \textbf{Unpacking compressed argument.} An author writing for trained listeners omits intermediate steps and relies on theorems argued questions earlier, packing whole inferences into a single clause introduced by \emph{enim} (``for'').
\item \textbf{Holding a distinction steady across a page.} A term is split in line three---self-evident in itself versus to us, prior in time versus in nature---and the whole argument collapses if the reader lets the two senses blur by line ten.
\item \textbf{Reading citations in context.} A quoted authority often functions as a shared premise or a disputed objection, never as a bare proof-stopper.
\end{itemize}

Our six questions---\textbf{Claim}, \textbf{Meaning}, \textbf{Warrant}, \textbf{Burden}, \textbf{Consequence}, and \textbf{Denial}---along with eight kinds of warrant and five kinds of denial, exist to train those three habits without requiring a medieval curriculum. Keep one rule inviolable: \textbf{these questions and categories are our modern reading devices, never the apparatus of Aristotle or Aquinas.} Aristotle never asked a ``Burden question,'' and Aquinas never sorted premises under our eight headings. Where a modern word of ours touches a classical term, the difference is as important as the overlap. Our ``metaphysical principle'' names any substantive starting claim about reality offered in an argument, contested or not; Aristotle's \emph{arche} is immediate and necessary. Our ``Warrant'' sorts whatever support a passage actually gives, including revelation or bare citation; Aristotle's six conditions define what belongs in a finished demonstrative science.

\subsection{Worked Passage 1: Validity Is Cheap, Demonstration Is Not}

Our first exhibit is an explicit reconstruction in the first syllogistic figure, which Aristotle calls the most scientific figure because it is the vehicle of mathematical proof and yields the reasoned fact (\emph{Posterior Analytics} I.14, 79a17--23). Compare two inferences built from the same three terms:

\begin{quote}
\emph{First inference:} Everything that changes has an unrealized capacity. This thing changes. Therefore, this thing has an unrealized capacity.
\end{quote}

Now swap the second premise and the conclusion:

\begin{quote}
\emph{Second inference:} Everything that changes has an unrealized capacity. This thing has an unrealized capacity. Therefore, this thing changes.
\end{quote}

Here symbolic notation earns its place for one specific purpose: displaying the bare logical form exposes a formally invalid inference---affirming the predicate instead of the subject---without relying on any theory of change or capacity. Let \linline{C} stand for ``thing that changes,'' \linline{U} for ``thing with an unrealized capacity,'' and \linline{a} for ``this thing.'' The first argument has the valid pattern:
\ldisplay{\logicop{All}\lmed C\lmed \logicop{are}\lmed U;\lquad a\lmed \logicop{is}\lmed C;\lquad \logicop{therefore}\lmed a\lmed \logicop{is}\lmed U.}
The second reverses the minor premise and conclusion into an invalid pattern:
\ldisplay{\logicop{All}\lmed C\lmed \logicop{are}\lmed U;\lquad a\lmed \logicop{is}\lmed U;\lquad \logicop{therefore}\lmed a\lmed \logicop{is}\lmed C.}

One homely counterexample kills the second pattern immediately: ``All burning logs are dry wood; this unlit log stacked by the hearth is dry wood; therefore, this unlit log is burning.'' Both premises are true, yet the conclusion is false. The second inference fails regardless of what act and potency turn out to be.

Now inspect the \emph{valid} first inference, granting for the moment that both premises are true. Is it an Aristotelian demonstration? \textbf{No.} It fails specific conditions of \emph{Posterior Analytics} I.2 (71b19--22) and I.4 (73b27--74a3). Nothing here shows that the first premise is \emph{primary} and \emph{immediate} rather than a definition or a derivative thesis. The second premise (``this thing changes'') states a contingent particular fact rather than a necessary universal truth. Most decisively, the middle term (\emph{changes}) is not the \emph{cause} of the thing's unrealized capacity: a log does not possess the capacity to burn because it is burning; the capacity is a precondition of burning, not its effect.

The outcome is modest and honest. This exhibit settles no metaphysical question about act, potency, or change, and says so plainly. It proves only that formal validity is cheap and never by itself constitutes a demonstration.

\subsection{Worked Passage 2: Showing That and Showing Why on One Conclusion}

Our second exhibit comes from Aristotle's \emph{Posterior Analytics} I.13, 78a30--38, in G.\ R.\ G.\ Mure's Oxford translation:

\begin{quote}
Thus you might prove as follows that the planets are near because they do not twinkle: let C be the planets, B not twinkling, A proximity. Then B is predicable of C; for the planets do not twinkle. But A is also predicable of B, since that which does not twinkle is near---we must take this truth as having been reached by induction or sense-perception. Therefore A is a necessary predicate of C; so that we have demonstrated that the planets are near. This syllogism, then, proves not the reasoned fact but only the fact; since they are not near because they do not twinkle, but, because they are near, do not twinkle.
\end{quote}

In the syllogism Aristotle quotes, the conclusion is \emph{that the planets are near}, and the middle term linking planets to proximity is \textbf{not twinkling}. On Aristotle's astronomical premise (that among heavenly lights, all and only nearby bodies shine without twinkling), both premises are true and the inference is valid. It achieves a genuine \emph{demonstratio quia}: starting from an observable effect better known to our eyes, it establishes \emph{that} the planets are near. Yet it fails to show \emph{why} they are near, because steady light does not cause closeness.

Now put the cause in the middle position:

\begin{quote}
Whatever is near does not twinkle. The planets are near. Therefore, the planets do not twinkle.
\end{quote}

Here the conclusion is \emph{that the planets do not twinkle}, and the middle term is \textbf{proximity}. Because proximity is the actual physical cause explaining why the planets shine steadily rather than shimmering like distant stars, this reversed syllogism is a \emph{demonstratio propter quid}---a demonstration of the reasoned fact.

Both arguments share the same valid first-figure form, the same three terms, and (on Aristotle's optical assumption) true premises. Validity cannot separate them. Everything turns on which term stands in the middle: the visible effect better known to us, or the real cause prior in nature.

\subsection{Exercises}

\subsubsection*{Exercise 0.1}
Right after the planets and twinkling, Aristotle gives a second illustration using the phases of the moon (\emph{Posterior Analytics} I.13, 78b3--10). Reconstructed in plain English:

\begin{itemize}
\item \textbf{(a)} Whatever waxes through crescent and half-circle phases is spherical. The moon waxes through crescent and half-circle phases. Therefore, the moon is spherical.
\item \textbf{(b)} Whatever is spherical waxes through crescent and half-circle phases when lit from one side. The moon is spherical (and lit from one side by the sun). Therefore, the moon waxes through crescent and half-circle phases.
\end{itemize}

Which argument is a demonstration \emph{quia}, and which is a demonstration \emph{propter quid}? Name the middle term in each, and explain why the two differ despite sharing valid form and true premises.

\subsubsection*{Exercise 0.2}
Consider this valid argument with true premises:

\begin{quote}
Every liquid that turns blue litmus paper red dissolves chalk. Vinegar turns blue litmus paper red. Therefore, vinegar dissolves chalk.
\end{quote}

Is this an Aristotelian demonstration \emph{propter quid}? Name at least two conditions from \emph{Posterior Analytics} I.2 (71b19--22) that it fails, and say what middle term would be required to turn it into one.

\subsubsection*{Exercise 0.3}
Read these two excerpts from \emph{Summa Theologiae} I, q.\,3, a.\,1 (``Whether God Is a Body?''), in the 1920 English Dominican translation. First, Objection 3:

\begin{quote}
Obj.\ 3: Further, whatever has corporeal parts is a body. Now Scripture attributes corporeal parts to God. ``Hast thou an arm like God?'' (Job 40:4); and ``The eyes of the Lord are upon the just'' (Ps.\ 33:16); and ``The right hand of the Lord hath wrought strength'' (Ps.\ 117:16). Therefore God is a body.
\end{quote}

Second, the opening of the \emph{corpus} and Reply to Objection 3:

\begin{quote}
\emph{I answer that,} It is absolutely true that God is not a body\dots\ \emph{Reply Obj.\ 3:} Corporeal parts are attributed to God in Scripture on account of His actions, and this is owing to a certain parallel. For instance the act of the eye is to see; hence the eye attributed to God signifies His power of seeing intellectually, not sensibly; and so on with the other parts.
\end{quote}

What error would a reader commit by quoting Objection 3 as Aquinas's view of Scripture? Explain what role the biblical quotations play in the \emph{disputatio} structure and how Reply 3 handles them.

\subsection{Solutions}

\subsubsection*{Solution to Exercise 0.1}
Argument (a) is a demonstration \emph{quia}, and argument (b) is a demonstration \emph{propter quid}. In (a), the middle term is \textbf{waxing through crescent and half-circle phases}: an observable effect better known to our sight that establishes \emph{that} the moon is spherical without explaining \emph{why}, since phases do not cause sphericity. In (b), the middle term is \textbf{being spherical} (under one-sided illumination), which names the geometric cause producing those phases. A common error is to dismiss (a) as invalid or merely inductive; on Aristotle's view, when effect and cause reciprocate, a \emph{quia} syllogism genuinely proves the fact, while stopping short of the reasoned fact.

\subsubsection*{Solution to Exercise 0.2}
No. Although valid and true, the argument fails two conditions of \emph{Posterior Analytics} I.2. First, its middle term---\emph{turning blue litmus paper red}---is not the \textbf{cause} of dissolving chalk; both color change and carbonate dissolution are co-effects of a shared chemical cause, namely acidity. Second, its major premise is neither \textbf{primary} nor \textbf{immediate}, since it rests on deeper facts about aqueous acids rather than standing as an undemonstrated starting point. At best, litmus color acts as a diagnostic sign; to approach a demonstration \emph{propter quid}, the middle term must name the actual chemical cause (being an aqueous acid reactive with calcium carbonate).

\subsubsection*{Solution to Exercise 0.3}
Treating Objection 3 as Aquinas's own teaching misreads the form of a scholastic \emph{disputatio}. Objections voice the strongest arguments against the thesis Aquinas will defend, frequently quoting respected authorities. Because Aquinas receives canonical Scripture as authoritative, he does not reject Job or the Psalms; after determining in the \emph{corpus} that God cannot be a body, he distinguishes in Reply 3 how Scripture speaks, attributing bodily parts ``on account of His actions'' by a parallel (``eye'' signifying intellectual sight, not a physical organ). Inside the article, the citations in Objection 3 function as shared authoritative texts posing a real interpretive difficulty to be resolved, never as a bare proof settling the question.
